\documentclass[12pt]{article}
\usepackage[utf8]{inputenc}
\usepackage[T1]{fontenc}
\usepackage[margin=1in]{geometry}
\usepackage{setspace}
\usepackage{xcolor}
\usepackage[colorlinks=true, urlcolor=blue!60!black, linkcolor=blue!60!black]{hyperref}
\setlength{\parindent}{2em}
\setlength{\parskip}{1.25em}
\renewcommand{\baselinestretch}{1.0}
\title{\textit{Zivotofsky v. Kerry}, 576 U.S. 1 (2015)}
\author{Case Brief}
\date{}

\begin{document}
\maketitle

\section{Case Metadata}
\begin{description}
  \item[Full Case Name:] \textit{Menachem Binyamin Zivotofsky, by His Parents and Guardians, Ari Z. and Naomi Siegman Zivotofsky v. John Kerry, Secretary of State}
  \item[Citation:] 576 U.S. 1 (2015) (Docket No. 13-628)
  \item[Argued / Decided:] November 3, 2014 / June 8, 2015
  \item[Vote:] 6--3 (5--4 as to the Court's Reception Clause rationale)
  \item[Author:] Justice Anthony M. Kennedy (joined by Justices Ginsburg, Breyer, Sotomayor, and Kagan)
  \item[Concurrences / Dissents:] Justice Breyer (concurring); Justice Thomas (concurring in the judgment in part and dissenting in part); Chief Justice Roberts (dissenting, joined by Justice Alito); Justice Scalia (dissenting, joined by Chief Justice Roberts and Justice Alito)
  \item[Chapter / Topic:] Separation of Powers and the Executive --- Presidential Recognition Power, Foreign Affairs, and \textit{Youngstown} Category Three
  \item[Primary Source URL:] \url{https://supreme.justia.com/cases/federal/us/576/1/}
\end{description}

\section{Facts and Procedural History}
Since President Harry S. Truman formally recognized the State of Israel on May 14, 1948, every subsequent President through the time of this litigation maintained a consistent Executive Branch policy declining to recognize any nation's sovereignty over the city of Jerusalem, insisting that the city's final status should be resolved through negotiations rather than unilateral action. Reflecting this policy, the State Department's Foreign Affairs Manual (7 FAM \S~1383) generally instructed consular officials to record a U.S. citizen's place of birth on a passport as the country having present sovereignty over the birthplace, but for citizens born in Jerusalem, officials were directed to record only ``Jerusalem'' and not ``Israel.''

In 2002, Congress enacted the Foreign Relations Authorization Act, Fiscal Year 2003, 116 Stat. 1350. Section 214, titled ``United States Policy with Respect to Jerusalem as the Capital of Israel,'' sought to counter Executive Branch neutrality. Specifically, \S~214(d) provided that ``[f]or purposes of the registration of birth, certification of nationality, or issuance of a passport of a United States citizen born in the city of Jerusalem, the Secretary shall, upon the request of the citizen or the citizen's legal guardian, record the place of birth as Israel.'' Upon signing the Act, President George W. Bush issued a signing statement declaring that \S~214, if construed as mandatory rather than advisory, would impermissibly interfere with the President's constitutional authority to speak for the Nation in foreign affairs and determine recognition, and affirming that U.S. policy regarding Jerusalem had not changed.

In 2002, Menachem Binyamin Zivotofsky was born in Jerusalem to U.S. citizen parents. In December 2002, his mother applied at the U.S. Embassy in Tel Aviv for a U.S. passport and a consular report of birth abroad, requesting that his birthplace be listed as ``Jerusalem, Israel.'' Embassy officials refused, listing only ``Jerusalem'' pursuant to State Department policy. Zivotofsky's parents sued the Secretary of State in the U.S. District Court for the District of Columbia to enforce \S~214(d). The District Court initially dismissed the suit for lack of standing and as a nonjusticiable political question. After the D.C. Circuit reversed on standing, 444 F.3d 614 (2006), the District Court again dismissed on political question grounds, and the D.C. Circuit affirmed, 571 F.3d 1227 (2009). In \textit{Zivotofsky v. Clinton}, 566 U.S. 189 (2012) (\textit{Zivotofsky I}), the Supreme Court vacated and remanded, holding that the case was justiciable because the Judiciary was asked only to decide the constitutionality of a federal statute, not the sovereign status of Jerusalem. On remand, the D.C. Circuit held \S~214(d) unconstitutional as an impermissible infringement on the President's exclusive recognition power, 725 F.3d 197 (2013). The Supreme Court granted certiorari.

\section{Issue(s) Presented}
\begin{enumerate}
  \item Does the Constitution vest the President with the exclusive power to grant or withhold formal recognition of a foreign sovereign and its territorial bounds?
  \item If so, does \S~214(d) of the Foreign Relations Authorization Act, Fiscal Year 2003---which commands the Secretary of State, upon request, to record ``Israel'' as the place of birth on the passport of a U.S. citizen born in Jerusalem---unconstitutionally infringe upon that exclusive presidential power?
\end{enumerate}

\section{Holding}
Yes to both questions. Affirming the D.C. Circuit, the Supreme Court held that (1) the President possesses the exclusive constitutional power to grant formal recognition to a foreign sovereign and its territorial bounds, and (2) \S~214(d) is unconstitutional because Congress may not command the President or the Secretary of State to issue a formal statement in an official Executive Branch document (a passport) that directly contradicts the President's prior recognition determination. (Because Zivotofsky waived any separate argument distinguishing consular reports of birth abroad from passports, the Court addressed only passports.)

\section{Reasoning}
Writing for a five-Justice majority, Justice Kennedy anchored the analysis in Justice Jackson's tripartite framework from \textit{Youngstown Sheet \& Tube Co. v. Sawyer}, 343 U.S. 579, 635--638 (1952) (Jackson, J., concurring). Because the President's refusal to enforce \S~214(d) was directly incompatible with the expressed will of Congress, the Executive operated in \textit{Youngstown} Category Three---where presidential power is ``at its lowest ebb'' and can prevail only if the asserted authority is both ``exclusive'' and ``conclusive,'' disabling Congress from acting upon the subject.

\textbf{1. Text, Structure, and Functionalism:} Although the Constitution does not use the word ``recognition,'' Justice Kennedy inferred the recognition power from Article II's text and structure, centering on the Reception Clause, Art. II, \S~3 (directing that the President ``shall receive Ambassadors and other public Ministers''). Relying on founding-era international law jurists (Vattel, Bynkershoek, and Grotius) and Alexander Hamilton's post-1793 shift in \textit{Pacificus No. 1}, the Court reasoned that receiving an ambassador was tantamount to acknowledging the sending state's sovereignty. This inference was reinforced by Article II's Treaty and Appointments Clauses (Art. II, \S~2, cl. 2): while Senate consent is required to ratify treaties or confirm ambassadors, only the President can initiate diplomatic channels, negotiate treaties, or nominate envoys. The Court declined to decide whether Article II's Vesting Clause supplied additional support. Functionally, recognition is a subject on which the Nation must ``speak . . . with one voice'' (\textit{American Ins. Assn. v. Garamendi}, 539 U.S. 396, 424 (2003)). Only the unitary Executive possesses the ``[d]ecision, activity, secrecy, and dispatch'' (\textit{The Federalist No. 70}) required to conduct delicate diplomatic contacts and issue unequivocal recognition determinations under international law.

\textbf{2. Precedent and the Narrowing of \textit{Curtiss-Wright}:} Reviewing prior decisions---including \textit{Williams v. Suffolk Ins. Co.}, 13 Pet. 415 (1839), \textit{United States v. Belmont}, 301 U.S. 324 (1937), \textit{United States v. Pink}, 315 U.S. 203 (1942), and \textit{Banco Nacional de Cuba v. Sabbatino}, 376 U.S. 398 (1964)---the Court found consistent judicial endorsement of exclusive executive control over recognition. Crucially, however, Justice Kennedy expressly rejected the Secretary of State's invitation to embrace the sweeping ``sole organ'' dicta of \textit{United States v. Curtiss-Wright Export Corp.}, 299 U.S. 304, 320 (1936). The Court stressed that \textit{Curtiss-Wright} involved congressionally authorized action, not unilateral defiance of a statute, and emphasized that Congress retains substantial lawmaking authority over foreign commerce, war, naturalization, appropriations, and passports.

\textbf{3. Historical Practice:} Invoking \textit{NLRB v. Noel Canning}, 573 U.S. 513 (2014), the Court surveyed historical episodes from President Washington's unilateral reception of Citizen Genet (1793) through South American independence (1818--1822), Texas (1836--1837), Haiti and Liberia (1861--1862), Cuba (1898), and President Carter's recognition of the People's Republic of China and the Taiwan Relations Act (1979). While acknowledging that history was ``not all on one side,'' the Court concluded that Congress had overwhelmingly acquiesced in presidential exclusivity over formal recognition and had never enacted a statute directly contradicting a President's formal recognition policy.

\textbf{4. Invalidity of \S~214(d):} Because formal recognition can be effected or undermined by official written declarations, the exclusive recognition power must include the President's authority to maintain his determination in his and his agents' official statements. A passport is a ``political document'' addressed to foreign powers (\textit{Urtetiqui v. D'Arcy}, 9 Pet. 692, 699 (1835)), and the place-of-birth field had long been treated as reflecting U.S. recognition policy. Moreover, \S~214's title, legislative history, and international reception demonstrated that Congress's purpose was to aggrandize its power at the expense of the Executive by forcing the Secretary of State to contradict presidential policy. While Congress retains broad authority over passports generally (e.g., \textit{Kent v. Dulles}, 357 U.S. 116 (1958)), it may not use that power to compel the Executive to contradict a formal recognition determination.

\section{Concurring and Dissenting Opinions}
\begin{itemize}
  \item \textbf{Justice Breyer (Concurring):} Maintained his dissent from \textit{Zivotofsky I} that the dispute presented a nonjusticiable political question, but joined the Court's opinion in full because precedent foreclosed dismissal on political question grounds.
  \item \textbf{Justice Thomas (Concurring in the Judgment in Part and Dissenting in Part):} Rejected the majority's Reception Clause rationale, arguing that listing ``Israel'' on a passport does not implicate international-law recognition because the United States has already recognized Israel as a state since 1948. Instead, Justice Thomas grounded his analysis in the ``residual foreign affairs power'' of Article II's Vesting Clause (contrasting Art. II, \S~1's omission of ``herein granted'' with Art. I, \S~1). Because passports have historically been executive travel documents addressed to foreign powers and Congress lacks an enumerated ``passport power'' under Article I or the ``proper'' prong of the Necessary and Proper Clause, \S~214(d) is unconstitutional as applied to passports. By contrast, consular reports of birth abroad are domestic citizenship documents tied to Congress's Naturalization Clause power; accordingly, Justice Thomas dissented from invalidating \S~214(d) as applied to birth reports.
  \item \textbf{Chief Justice Roberts (Dissenting, joined by Justice Alito):} Emphasized that for the first 225 years of the Republic, the Court had never upheld a President's direct defiance of an Act of Congress in foreign affairs. He argued that the Reception Clause is phrased as a ministerial duty (``shall receive''), not an exclusive grant of power, citing Hamilton's \textit{The Federalist No. 69}, and that conflicting precedent and history failed to establish ``conclusive and preclusive'' presidential authority. Even assuming exclusive recognition power, \S~214(d) did not implicate it: an optional birthplace accommodation on a passport is not an act of formal recognition, and giving dispositive weight to foreign observers' misperceptions subjects duly enacted laws to an ``international heckler's veto.''
  \item \textbf{Justice Scalia (Dissenting, joined by Chief Justice Roberts and Justice Alito):} Grounded \S~214(d) squarely in Congress's Naturalization Clause and Necessary and Proper Clause powers to issue citizenship-verifying documents. He drew a sharp distinction between formal recognition---a binding legal act creating international-law estoppel---and a geographic accommodation in a passport, noting the State Department's analogous accommodations for Belfast, pre-1948 Palestine, and Taiwan. He mocked the majority's extension of recognition to official ``statements'' as ``a leap worthy of the Mad Hatter,'' warned that functionalist ``one voice'' reasoning systematically favors the unitary President over the plural Congress, and criticized Justice Thomas's residual-powers theory as producing ``a presidency more reminiscent of George III than George Washington.''
\end{itemize}

\section{Key Quotations}
\begin{quote}
``The President does have a unique role in communicating with foreign governments, as then-Congressman John Marshall acknowledged. \dots\ But whether the realm is foreign or domestic, it is still the Legislative Branch, not the Executive Branch, that makes the law.'' \\
--- Justice Kennedy, for the Court, 576 U.S. at 21.
\end{quote}

\begin{quote}
``The Executive is not free from the ordinary controls and checks of Congress merely because foreign affairs are at issue. \dots\ It is not for the President alone to determine the whole content of the Nation's foreign policy.'' \\
--- Justice Kennedy, for the Court, 576 U.S. at 21.
\end{quote}

\begin{quote}
``This case is confined solely to the exclusive power of the President to control recognition determinations, including formal statements by the Executive Branch acknowledging the legitimacy of a state or government and its territorial bounds. Congress cannot command the President to contradict an earlier recognition determination in the issuance of passports.'' \\
--- Justice Kennedy, for the Court, 576 U.S. at 32.
\end{quote}

\begin{quote}
``Today's decision is a first: Never before has this Court accepted a President's direct defiance of an Act of Congress in the field of foreign affairs.'' \\
--- Chief Justice Roberts, dissenting, 576 U.S. at 61.
\end{quote}

\begin{quote}
``Recognition is a type of legal act, not a type of statement. It is a leap worthy of the Mad Hatter to go from exclusive authority over making legal commitments about sovereignty to exclusive authority over making statements or issuing documents about national borders.'' \\
--- Justice Scalia, dissenting, 576 U.S. at 75.
\end{quote}

\section{Doctrinal Significance and Modern Legacy (Through 2026)}
\textit{Zivotofsky v. Kerry} is a landmark in modern separation-of-powers jurisprudence for two competing reasons. First, it represents the first time in Supreme Court history that a President prevailed in \textit{Youngstown} Category Three in the foreign-affairs arena, establishing that the Reception Clause and related Article II provisions carve out a ``conclusive and preclusive'' executive sphere that disables Congress from contradicting presidential recognition decisions. Second, the Court simultaneously delivered the sharpest modern repudiation of \textit{Curtiss-Wright}'s ``sole organ'' dicta, reaffirming that Congress remains the primary lawmaker in foreign affairs as well as domestic affairs.

The decision's practical epilogue underscored the real-world force of exclusive executive recognition: in December 2017, President Donald J. Trump unilaterally recognized Jerusalem as Israel's capital, and in October 2020 the State Department revised its passport policy and presented Menachem Zivotofsky himself with the first U.S. passport listing ``Israel'' as the place of birth for a Jerusalem-born citizen (\url{https://www.timesofisrael.com/us-gives-first-passport-with-israel-birthplace-to-jerusalem-born-teen/}).

Through 2026, both halves of \textit{Zivotofsky} have proved foundational at the Supreme Court:
\begin{itemize}
  \item \textbf{Expansion of ``Conclusive and Preclusive'' Authority in \textit{Trump v. United States}, 603 U.S. 593 (2024):} Writing for the Court, Chief Justice Roberts---despite his dissent in \textit{Zivotofsky}---relied directly on \textit{Zivotofsky} (citing 576 U.S. at 32) as a canonical illustration of the President's ``conclusive and preclusive'' constitutional authority, reasoning that Congress cannot act on (including by criminalizing), and courts cannot examine, presidential conduct within that exclusive Article II sphere (\url{https://www.law.cornell.edu/supremecourt/text/23-939}).
  \item \textbf{Reaffirmation of Congressional Lawmaking Primacy in \textit{Learning Resources, Inc. v. Trump}, 607 U.S. \_\_\_ (2026):} On February 20, 2026, the Supreme Court held 6--3 that the International Emergency Economic Powers Act (IEEPA) does not authorize the President to unilaterally impose tariffs, underscoring that the taxing and tariff powers belong to Congress under Article I rather than to an unbounded executive foreign-affairs prerogative (\url{https://supreme.justia.com/cases/federal/us/607/24-1287/}). Meanwhile, Justice Thomas's residual foreign-affairs Vesting Clause concurrence in \textit{Zivotofsky} (576 U.S. at 36) continues to be cited in 2025--2026 litigation over executive authority in trade, diplomacy, and immigration.
\end{itemize}

\section{Pedagogical Critiques and Exam Hypotheticals}
\textbf{Pedagogical Critiques:}
\begin{itemize}
  \item \textit{Functionalism vs. Formalism:} The majority blends textual inference (the Reception Clause) with pragmatic functionalism (``speaking with one voice''), whereas both Justice Scalia and Justice Thomas critique the opinion from formalist angles---demonstrating that ``originalism'' itself fractures sharply in foreign affairs between pro-congressional enumerated-powers formalism (Scalia) and pro-executive Vesting Clause residualism (Thomas).
  \item \textit{The ``Statement'' vs. ``Legal Act'' Slippage:} Students should scrutinize the majority's admission that \S~214(d) ``would not itself constitute a formal act of recognition'' alongside its claim that the President's exclusive power ``extends no further than his formal recognition determination.'' How far does the ancillary right to control ``official executive statements'' reach before swallowing Congress's enumerated foreign-commerce and naturalization powers?
  \item \textit{Asymmetric \textit{Noel Canning} Historical Gloss:} Chief Justice Roberts highlights a crucial doctrinal tension: historical congressional acquiescence normally supports presidential power in \textit{Youngstown} Category Two (congressional silence), not Category Three (direct statutory defiance).
\end{itemize}

\textbf{Exam Hypothetical (and ``Opposite Side'' Contrast):} \\
\textit{Hypothetical:} Suppose the President formally recognizes the government of Country X as the sole legitimate sovereign over Disputed Island Y, refusing to recognize a breakaway republic on Island Y. Disagreeing with the President, Congress enacts the ``Island Y Trade and Travel Act,'' which contains two provisions: (1) Section A directs U.S. Customs and Border Protection to classify all goods imported from Island Y as ``Products of the Independent Republic of Island Y'' on official customs entry invoices and to apply a 0\% preferential tariff rate rather than Country X's 25\% tariff rate; and (2) Section B prohibits the expenditure of any federal funds to operate a U.S. Embassy in Country X's capital unless the President withdraws recognition of Country X's claim to Island Y. The President signs the bill with a signing statement refusing to comply with Section A and challenging Section B.

\textit{Analysis \& ``Opposite Side'' Contrast:}
\begin{itemize}
  \item \textit{As to Section A (Customs Invoices and Tariffs):} Under \textit{Zivotofsky}'s majority opinion, the President will argue that compelling Customs officials to print ``Independent Republic of Island Y'' on official government documents forces the Executive to contradict its formal recognition determination in an official communication affecting foreign sovereigns. \textbf{Opposite Side Contrast:} Congress (echoing Justice Scalia's dissent and Kennedy's own limits on \textit{Curtiss-Wright}) has a much stronger counterargument here than in \textit{Zivotofsky}: regulating tariffs and country-of-origin rules lies at the core of Congress's express Foreign Commerce Clause and Taxing Clause powers (Art. I, \S~8, cls. 1, 3)---a power reinforced by \textit{Learning Resources, Inc. v. Trump} (2026)---and Congress has historically enacted trade and statutory rules treating unrecognized entities (such as Taiwan under the Taiwan Relations Act) as distinct for commercial purposes without altering formal diplomatic recognition. A court might sever the 0\% tariff rate (valid under Article I) from any mandatory sovereignty label on official diplomatic-facing documents, or uphold Section A entirely as a commercial classification rather than an act of recognition.
  \item \textit{As to Section B (Appropriations Condition):} \textbf{Opposite Side Contrast:} Even under \textit{Zivotofsky}'s majority opinion, Justice Kennedy explicitly acknowledged that Congress retains its Appropriations Clause power (Art. I, \S~9, cl. 7) and that ``[t]he President, furthermore, could not build an American Embassy abroad without congressional appropriation of the necessary funds.'' Thus, refusing to fund an embassy is squarely within Congress's power of the purse. However, the President's counterargument would invoke the unconstitutional-conditions principle in separation of powers: Congress may not use an explicit condition tied directly to forcing a change in a formal recognition determination to coerce the President into surrendering an exclusive Article II prerogative.
\end{itemize}

\end{document}
